\section{\ref{sec:dofaalt}장. \nt{DOPA}의 다른 이론들}

확장된 그림들은 저마다 도파민의 직접 측정(반전점의 흩어짐, 불쾌한 자극과 새로운 자극에 대한 반응, 완만한 상승, 맛 교체에 대한 반응)에 기대고 있다. 그러나 그것을 설명하는 식들은 이론이며, 그 가운데 두 이론 사이에서는 실험 결과가 아직 서로 모순된다.

{\footnotesize
\setlength{\tabcolsep}{3pt}\renewcommand{\arraystretch}{1.15}
\begin{longtable}{>{\raggedright}p{0.5cm}>{\raggedright}p{4.0cm}>{\raggedright}p{5.6cm}>{\raggedright}p{2.3cm}>{\raggedright\arraybackslash}p{1.7cm}}
\toprule
No. & 명제 & 실험과 측정한 것 & 출처 & 상태 \\
\midrule\endfirsthead
\toprule
No. & 명제 & 실험과 측정한 것 & 출처 & 상태 \\
\midrule\endhead
1 & 비대칭 규칙~\eqref{eq:asymmetric}은 점~\eqref{eq:expectile}으로 수렴한다. $\kappa=1/2$이면 평균이고, 확률 $1/2$의 한 방울이면 $V=\kappa$이다(그림~\ref{fig:expectiles}) & 점~\eqref{eq:expectile}은 익스펙타일, 곧 양의 편차와 음의 편차에 서로 다른 가중치를 준 최소제곱 문제의 해다. 이 장의 예에 대한 유도는 본문에 있다. & \cite{NeweyPowell1987}; 이 장의 유도 & 이론 \\
2 & \nt{DOPA} 뉴런마다 반전점이 다르고, 반전점은 비대칭 순으로 정렬된다(명제~\ref{prop:expectile}) & 생쥐, 배쪽 피개 영역의 \nt{DOPA} 뉴런을 광유전학으로 표지, 크기가 다른 보상: 급증이 급감으로 바뀌는 점이 뉴런마다 다르다. 반응 비대칭이 큰 뉴런일수록 그 점이 높다. 무리 반응으로 보상 분포의 모양이 복원된다. 이것이 흩어짐의 주된 기능이라는 것은 가설이다. & \cite{Dabney2020} & 측정됨 \\
3 & 일부 \nt{DOPA} 뉴런은 불쾌한 자극에 급증으로 반응한다 & 원숭이: 어떤 \nt{DOPA} 뉴런은 보상 예고 신호에 흥분하고 눈에 부는 공기 예고 신호에 억제되며, 다른 뉴런은 둘 다에 흥분한다. 두 무리는 중뇌의 서로 다른 부위에 있다. & \cite{Matsumoto2009, BrombergMartin2010} & 측정됨 \\
4 & 오차의 크기 또는 새로움이 학습률을 정한다 & \nt{DOPA} 뉴런의 중요도 신호가 학습률을 바꾼다는 직접 실험은 인용된 연구에 없다. 리뷰는 “주의, 중요함” 신호의 역할을 제안한다. & \cite{BrombergMartin2010} & 가설 \\
5 & 위험 축을 위한 별도의 전선 & 생쥐: 선조체 꼬리로 가는 \nt{DOPA} 뉴런은 보상이 아니라 새롭고 강한 자극에 반응한다. 이 뉴런을 파괴하면 새로운 물체에 대한 회피가 약해진다. & \cite{Menegas2018} & 측정됨 \\
6 & 최적 행동 시간 $\tau_a^*=\sqrt{C/\bar R}$~\eqref{eq:vigor} & 합 $C/\tau_a+\bar R\tau_a$의 최소. 유도는 본문에 있다. 원래 모델에서 행동 속도는 평균 보상과 같은 시간의 가격이 정한다. & \cite{Niv2007}; 이 장의 유도 & 이론 \\
7 & 느린 도파민 수준은 $\bar R$를 싣고 행동 속도를 정한다(명제~\ref{prop:multiplex}) & 쥐, 측좌핵의 미세투석과 전압전류법: 도파민 수준이 분 단위로 보상 빈도와 행동 속도를 따라간다. 사람에서 L-DOPA는 반응 시간이 평균 보상 빈도에 의존하는 정도를 키운다. 느린 수준이 정확히 $\bar R$라는 것은 증명되지 않았다. & \cite{Hamid2016, Beierholm2013vigor} & 가설 \\
8 & 느린 수준은 두 원천에서 합쳐진다. 방출은 스파이크 발화율이 변하지 않아도 출력 말단에서 달라지지만, 발화율 자체에도 완만한 상승이 있다 & 쥐, 한 과제: 측좌핵의 방출은 변하는 보상 기대를 따라가지만, 식별된 배쪽 피개 영역 \nt{DOPA} 뉴런의 발화율은 그렇지 않다. 가상 복도의 생쥐(10행): 보상에 다가갈 때의 완만한 상승이 식별된 \nt{DOPA} 뉴런의 스파이크에도 있다. & \cite{Mohebi2019, Kim2020} & 측정됨 \\
9 & 동물이 먼 보상에 다가가는 동안 도파민이 완만하게 오른다 & 미로의 쥐, 선조체 전압전류법: 목표에 가까워지면서 도파민 농도가 몇 초에 걸쳐 오르고, 기울기는 거리와 보상 크기에 의존한다. & \cite{Howe2013ramp, Hamid2016} & 측정됨 \\
10 & 상승은 추정이 정확해질 때의 오차 $\delta$다~\eqref{eq:ramp}. 복도에서의 도약은 급증을 준다(그림~\ref{fig:ramp}) & 식과 수치 초당 $\delta=0.75\,V$는 본문의 유도. 실험: 가상 복도의 생쥐를 목표 쪽으로 순간 이동. 세포체의 스파이크, 세포체와 출력 말단의 칼슘, 선조체의 농도가 기대 $V$가 아니라 오차처럼 반응한다. 두 설명 중 어느 것이 모든 출력에서 옳은지는 논쟁 중이다. & \cite{Kim2020} & 측정됨 \\
11 & 뒤돌아보기: 도파민은 인과 연관의 척도~\eqref{eq:retro}를 알린다 & 생쥐, 이 이론과 오차~\eqref{eq:ctd}가 서로 다른 것을 예측하는 실험에서의 측좌핵 도파민 방출: 여러 실험에서 방출이 앞의 이론을 따랐다. & \cite{Jeong2022} & 가설 \\
12 & 학습 중 도파민 반응은 보상에서 예고 신호로 점차 옮겨 간다 & 생쥐, 학습 과정 동안 배쪽 선조체에서 \nt{DOPA} 뉴런 출력 말단의 신호: 급증이 보상 순간에서 예고 신호 순간으로 점차 옮겨 간다. \eqref{eq:ctd}에 의한 학습이 요구하는 그대로다. & \cite{Amo2022} & 측정됨 \\
13 & \nt{DOPA} 뉴런은 가치가 같아도 보상의 맛이 바뀌면 반응한다 & 쥐, 배쪽 피개 영역 뉴런 기록: 주스를 가치가 같은 다른 주스로 바꾸면 오차~\eqref{eq:ctd}가 0인데도 반응이 생긴다. & \cite{Takahashi2017} & 측정됨 \\
14 & 인위적 급증은 두 중립 자극 사이의 연관을 가르친다 & 쥐, 보상 없이 두 자극을 미리 짝짓는 실험: 첫째 자극에서 둘째 자극으로 넘어갈 때 \nt{DOPA} 뉴런을 광유전학으로 켜면 연관이 생기고, 끄면 학습이 방해된다. & \cite{Sharpe2017} & 측정됨 \\
15 & 특징별 오차 벡터~\eqref{eq:vectortd} & 특징 예측 오차 이론. 벡터 형태에 대한 직접 검증은 없다. & \cite{Gardner2018} & 가설 \\
16 & 한 과제에서 \nt{DOPA} 뉴런마다 서로 다른 양을 싣는다 & 가상 복도의 생쥐, \nt{DOPA} 뉴런의 2광자 영상: 어떤 뉴런은 보상과, 어떤 뉴런은 속도와 방향 전환과, 또 어떤 뉴런은 예고 신호와 선택의 정확도와 연관된다. & \cite{Engelhard2019} & 측정됨 \\
17 & 전선 대신 다발(명제~\ref{prop:bundle}) & 2--16행의 요약: 각 분할은 따로따로 측정되었다. 이것들이 모두 한 신호의 부분이라는 통합된 그림은 실험으로 확인되지 않았다. & --- & 가설 \\
\bottomrule
\end{longtable}}
